\documentclass[popl.tex]{subfiles}
\begin{document}

\clearpage\section{Method}\label{sec:method}

The key observation is that finite CFL--regular intersection can be treated as CFL reachability over an acyclic automaton. Given a CFG \(G\) and an acyclic NFA \(A\), we compute which nonterminals derive the labels of paths \(p\rightsquigarrow r\) in \(A\), then ask whether some accepting path from \(q_\alpha\) derives the start symbol. This is close in spirit to highly parallel CKY recognition: Brent and Goldschlager show that a string of length \(n\) can be parsed in \(\mathcal{O}(\log n)\) time on a unit-cost SIMDAG, or \(\mathcal{O}(\log^2 n)\) time on a log-cost SIMDAG, using \(\mathcal{O}(n^6)\) processors~\cite{brent1984parallel}. Our setting is different: the right-hand input is not one string but a finite automaton, and the construction below is meant to produce either a Boolean decision or, in a second pass, a compact witness object.

\begin{theorem}\label{thm:parallel_decision_complexity}
  For any CFG, \(G = \langle V, \Sigma, P, S\rangle\), and acyclic NFA (ANFA), \(A = \langle Q, \Sigma, \delta, q_\alpha: Q, F \subseteq Q\rangle\), there exists a decision procedure \(\Psi: \text{CFG} \rightarrow \text{ANFA} \rightarrow \mathbb{B}\) such that \(\Psi(G, A) \models [\mathcal{L}(G)\cap\mathcal{L}(A) \neq \varnothing]\) requiring an \(\mathcal{O}(\log^2|Q|)\) topological-ordering phase followed by at most \(\mathcal{O}(|Q||V|)\) parallel closure rounds using \(\mathcal{O}\big(|Q|^2|V|\big)\) parallel random access (PRAM) processors.
\end{theorem}

\begin{proof}[Proof]
  To prove nonemptiness, we must show there exists a path $q_\alpha \rightsquigarrow q_\omega$ in $A$ such that $q_\omega: F$ where $q_\alpha \rightsquigarrow q_\omega \vdash S$. At least one of two cases must hold for $w \in V$ to parse a given $p \rightsquigarrow r$ pair:

  \begin{enumerate}
    \item $p$ steps directly to $r$ in which case it suffices to check $\exists a.\big((p \overset{a}{\rightarrow} r)\in \delta \land (w \rightarrow a) \in P\big)$, or, \item there is some midpoint $q \in Q$, $p \rightsquigarrow q \rightsquigarrow r$ such that $\exists x, z.\big((w \rightarrow xz) \in P\land\overbrace{\underbrace{p \rightsquigarrow q}_x, \underbrace{q \rightsquigarrow r}_z}^w\big)$.
  \end{enumerate}

  \noindent This decomposition immediately suggests a dynamic programming solution. Let M be a matrix of type $E^{|Q|\times|Q|\times|V|}$ indexed by $Q$. Since we assumed $\delta$ is acyclic, there exists a topological sort of $\delta$ imposing a total order on $Q$ such that $M$ is strictly upper triangular (SUT). Note $Q$ can be ordered topologically in $\mathcal{O}(\log^2 |Q|)$ time~\cite{dekel1981parallel} using matrix multiplication. We initialize $M$ thusly:
  \begin{align}
    M_0[p, r, w] = \bigvee_{a\in\Sigma}\big[(p \overset{a}{\rightarrow} r)\in\delta \land (w\rightarrow a)\in P\big].
  \end{align}
  For Boolean charts \(X,Y\), we define
 \begin{align}
    [X\otimes Y][p,r,w] = \bigvee_{\mathclap{q\in Q,\;x,z\in V}} \big[X[p,q,x]\land Y[q,r,z]\land (w\rightarrow xz)\in P\big].
  \end{align}
 then iterate the following fixedpoint until convergence:
  \[
    T_0=M_0,\qquad T_{j+1}=T_j \lor (T_j\otimes T_j).
  \]
  Here \(T_j^2\) is only shorthand for \(T_j\otimes T_j\). No associativity of \(\otimes\) is used: the algorithm never forms an unparenthesized \(M^3\), and each update corresponds to one grammar-bracketed split \(w\rightarrow xz\).

  Let \(H_j[p,r,w]\) be the assertion that some word labeling a path \(p\rightsquigarrow r\) is derived from \(w\) by a CNF parse tree of height at most \(j\). We claim \(T_j[p,r,w]=H_j[p,r,w]\) for all \(j\). The case \(j=0\) is \(M_0\). The step follows because a height-\((j+1)\) tree is either already of height \(j\), or has root \(w\rightarrow xz\) with two height-\(\leq j\) subtrees whose yields label paths \(p\rightsquigarrow q\) and \(q\rightsquigarrow r\). This is exactly the displayed product. Since \(A\) is acyclic, every accepting path has fewer than \(|Q|\) arcs, hence every CNF parse over it has height \(<|Q|\). Thus, the recurrence reaches its least fixed point \(T_\star\) within the stated \(\mathcal{O}(|Q||V|)\) rounds, using one processor per \((p,r,w)\) entry per round. Finally,
  \begin{align}
    \Psi(G,A)= \bigvee_{\mathclap{q_\omega\in F}} T_\star[q_\alpha,q_\omega,S],
  \end{align}
  which is equivalent to \(\mathcal{L}(G)\cap\mathcal{L}(A)\neq\varnothing\). The test may be performed after each round, giving positive early termination.
\end{proof}

\begin{theorem}[Witness construction]\label{thm:witness_construction}
  Under the hypotheses of Theorem~\ref{thm:parallel_decision_complexity}, the
  saturated chart can be refined to a finite star-free generalized regular
  expression \(S_\cap\) satisfying
  \(\mathcal{L}(S_\cap)=\mathcal{L}(G)\cap\mathcal{L}(A)\). Consequently, when
  \(\Psi(G,A)=\bs\), a word
  \(\sigma\in\mathcal{L}(G)\cap\mathcal{L}(A)\) can be decoded by repeatedly
  choosing a token whose Brzozowski derivative is nonempty.
\end{theorem}

\begin{definition}[Brzozowski derivative]\label{def:brzozowski_derivative}
  Following Brzozowski~\cite{brzozowski1964derivatives}, let \(E\) be the star-free GREs
  \(e ::= \varnothing\mid\varepsilon\mid a\mid e\vee e\mid e\land e\mid e\cdot e\).
  Write \(\delta(e)\in\{\varnothing,\varepsilon\}\) for the nullable projection,
  where \(\delta(e)=\varepsilon\) iff \(\varepsilon\in\mathcal{L}(e)\).  For
  \(a,b\in\Sigma\), define:

  \vspace{-0.8cm}
  \begin{multicols}{2}
    \begin{eqnarray*}
      \phantom{--}\partial_a(&\hspace{-0.35cm} \varnothing \hspace{-0.35cm}&) = \varnothing                                           \\
      \phantom{--}\partial_a(&\hspace{-0.35cm} \varepsilon \hspace{-0.35cm}&) = \varnothing                                           \\
      \phantom{--}\partial_a(&\hspace{-0.35cm} b           \hspace{-0.35cm}&) = \begin{cases}\varepsilon &\text{ if } a = b\\ \varnothing &\text{ if } a \neq b \end{cases}\\
      \phantom{--}\partial_a(&\hspace{-0.35cm} x\cdot z    \hspace{-0.35cm}&) = (\partial_a x)\cdot z \vee \delta(x)\cdot\partial_a z \\
      \phantom{--}\partial_a(&\hspace{-0.35cm} x\vee  z    \hspace{-0.35cm}&) =  \partial_a x \vee  \partial_a z                       \\
      \phantom{--}\partial_a(&\hspace{-0.35cm} x\land z    \hspace{-0.35cm}&) =  \partial_a x \land \partial_a z
    \end{eqnarray*} \break\vspace{-0.45cm}
    \begin{eqnarray*}
      \delta(&\hspace{-0.35cm} \varnothing \hspace{-0.35cm}&) = \varnothing                                      \\
      \delta(&\hspace{-0.35cm} \varepsilon \hspace{-0.35cm}&) = \varepsilon                                      \\
      \delta(&\hspace{-0.35cm} a           \hspace{-0.35cm}&) = \varnothing\phantom{\begin{cases}\varepsilon\\\varnothing\end{cases}}\\
      \delta(&\hspace{-0.35cm} x\cdot z    \hspace{-0.35cm}&) = \delta(x) \land \delta(z)                        \\
      \delta(&\hspace{-0.35cm} x\vee  z    \hspace{-0.35cm}&) = \delta(x) \vee  \delta(z)                        \\
      \delta(&\hspace{-0.35cm} x\land z    \hspace{-0.35cm}&) = \delta(x) \land \delta(z)
    \end{eqnarray*}
  \end{multicols}
  Extend to strings by \(\partial_\varepsilon e=e\) and
  \(\partial_{a\sigma}e=\partial_\sigma(\partial_a e)\); then
  \(\mathcal{L}(\partial_\sigma e)=\sigma^{-1}\mathcal{L}(e)\).
\end{definition}

\begin{lemma}[Generation]\label{lem:generation}
  For any finite GRE \(e\) with \(\mathcal{L}(e)\neq\varnothing\), the following
  nondeterministic procedure terminates and returns a word in \(\mathcal{L}(e)\):
  \begin{equation}\label{eq:choose}
    \texttt{choose}(e)=
    \begin{cases}
      \varepsilon & \delta(e)=\varepsilon,\\
      a\cdot\texttt{choose}(\partial_a e) &
      a\in\Sigma,\ \mathcal{L}(\partial_a e)\neq\varnothing .
    \end{cases}
  \end{equation}
\end{lemma}

\begin{proof}
  If \(\delta(e)=\varepsilon\), the empty word is a valid witness.  Otherwise every
  word in the nonempty finite language \(\mathcal{L}(e)\) has the form
  \(a\sigma\), so \(\mathcal{L}(\partial_a e)\neq\varnothing\) for at least one
  \(a\).  Taking such an \(a\) preserves nonemptiness and removes one leading
  symbol from every selected witness.  Finiteness rules out an infinite descent.
  If the emitted word is \(\tau\), the derivative law gives
  \(\varepsilon\in\mathcal{L}(\partial_\tau e)\), hence
  \(\tau\in\mathcal{L}(e)\).
\end{proof}

\begin{proof}
  Let \(E\) be star-free regexes
  \(e ::= \varnothing\mid\varepsilon\mid a\mid e\vee e\mid e\cdot e\), modulo
  language equivalence, with provenance semiring operations
  \(e\oplus f\coloneqq e\vee f\), \(e\otimes f\coloneqq e\cdot f\),
  \(0_E\coloneqq\varnothing\), and \(1_E\coloneqq\varepsilon\).  Keeping the
  Boolean support \(T_\star\) fixed, set unsupported cells to \(0_E\) and define
  \[
  \begin{gathered}
    R_0[p,r,w]=
    \bigoplus_{\mathclap{\substack{a\in\Sigma\\(p\overset{a}{\rightarrow}r)\in\delta\\(w\rightarrow a)\in P}}} a,\qquad
    R_{j+1}=R_j\oplus(R_j\otimes R_j),\\
    [X\otimes Y][p,r,w]=
    \bigoplus_{\mathclap{\substack{q\in Q,\;x,z\in V\\(w\rightarrow xz)\in P}}}
    X[p,q,x]\otimes Y[q,r,z],
  \end{gathered}
  \]
  where empty sums are \(0_E\).  This is the Boolean closure lifted to regex
  provenance: terminals record arcs, and products record a grammar rule plus a
  midpoint.  Topological ordering of \(A\) makes every nonzero product factor
  into strict subpaths, so the shared expression graph is acyclic and star-free.

  Induction on parse height shows \(R_j[p,r,w]\) denotes exactly the words
  labeling paths \(p\overset{\sigma}{\rightsquigarrow}r\) derivable from \(w\)
  by height-\(\leq j\) trees.  Acyclicity bounds the necessary height, hence
  \(R_j\) reaches \(R_\star\); with
  \(S_\cap=\bigoplus_{q_\omega\in F} R_\star[q_\alpha,q_\omega,S]\), the invariant gives
  \(\mathcal{L}(S_\cap)=\mathcal{L}(G)\cap\mathcal{L}(A)\).

  If \(\Psi(G,A)=\bs\), then \(S_\cap\) is nonempty.  Since it is finite and
  star-free, Lemma~\ref{lem:generation} applies and
  \(\texttt{choose}(S_\cap)\in\mathcal{L}(S_\cap)\), yielding the desired
  witness.
\end{proof}
\ifSubfilesClassLoaded{\par\clearpage\bibliography{../bib/acmart}}{}
\end{document}
